% ----------------------------------------------------------------------------------
% Chapter 5: Experiments and Results
% ----------------------------------------------------------------------------------
% Structure the results logically. Start with quantitative data, use figures, and avoid emotive adjectives.
\label{ch:experiments}

\section{Experimental Setup}
The tests were conducted in both simulated (Gazebo) and physical environments...

\section{Mapping and Navigation Performance}
% Metric data on localization accuracy and path execution time.

\section{Manipulation Precision}
% Success rate of pick-and-place routines.

\subsection{Reachability as the Binding Constraint}
The dominant finding of the manipulation experiments is that workspace reachability,
rather than planning or perception, governs grasp feasibility for this platform. The
nominal $0.28$\,m specification of the myCobot~280 describes the maximum reach of the
arm extended horizontally; the effective reach for a \emph{top-down} grasp is
substantially shorter, because orienting the gripper vertically requires the wrist to
fold, consuming reach. From the measured link geometry the elbow-to-wrist span is
approximately $0.226$\,m, which is insufficient to keep the wrist above a target
located $\approx 0.24$\,m forward; consequently a purely top-down grasp at that
distance admits no inverse-kinematics solution. Empirical reach maps, generated by
sweeping candidate poses through the \texttt{/compute\_ik} service and recording
success, confirmed this analysis.

\subsection{Orientation-Dependent Reachability}
A second result is that reachability is a function of the full target orientation,
including the redundant wrist roll, and not of position alone. For an identical target
position, one top-down wrist roll over-twisted the wrist at extension and yielded no
solution, whereas a roll rotated by $90^{\circ}$ was solvable across the workspace.
This sensitivity is significant for grasp planning: fixing the tool orientation to a
single value can make the usable workspace appear materially smaller than it is, and
reach must therefore be evaluated over a family of orientations. The interaction of
this short top-down reach with the mobile base's forward overhang and the
forward-mounted camera (Chapter~4) defines a narrow feasible window jointly governed by
arm mounting, base stop distance, and workspace layout---a constraint specific to
mobile manipulation and absent from fixed-base systems.

\subsection{Workspace Reachability Map}
The reachability analysis was quantified by sweeping candidate end-effector poses across
a grid of positions and orientations and recording inverse-kinematics success via the
\texttt{/compute\_ik} service. In a representative $2025$-pose sweep, the top-down
orientation was reachable at far fewer grid points than an angled or side approach
($194$ reachable poses top-down versus $290$ side-grasp), confirming that the top-down
orientation is the binding constraint and that a side or tilted approach materially
enlarges the usable workspace. An independent $750$-pose sweep recorded $25.9\%$ overall
reachability, with the reachable envelope contracting sharply beyond $x\approx0.18$\,m.
These maps were produced by an automated sweep tool and are treated as primary
experimental data, providing the empirical basis for the chosen grasp distance and
orientation.

\subsection{Grasp Execution in Simulation}
Stable grasp execution in Gazebo required addressing rigid-body contact behaviour. A
position-controlled gripper commanded into a rigid surface produces large separating
impulses that, with a comparatively light and unbraked base, destabilise the entire
platform. This was mitigated by (i) terminating the descent clear of supporting
surfaces and registering them as planning-scene collision objects, (ii) constraining
the base during manipulation, and (iii) adopting the established simulation convention
of rigidly attaching the object to the gripper at the grasp instant rather than relying
on simulated friction closure. These measures are characteristic of pick-and-place in
Gazebo and are reported here as part of the methodology rather than as a limitation of
the approach. With the tool-frame calibration of Chapter~4 applied (lateral offset,
tightened orientation tolerance, and closing clearance) and the surface contact model
softened, the simulated manipulator executes a centred, vertical top-down grasp that
lifts, holds, and replaces the target repeatably.

\section{System Integration and Logistics Simulation}
% The holistic test (multi-stage assembly logistics sequence).
